\begin{abstract}
\noindent Case counts are often taken to set the precision of reproduction-number estimates, yet the uncertainty that additional cases remove depends on the inferential target. In a one-generation branching model with negative binomial offspring and binomial reporting, we derive exact mean squared errors for three targets: the current reproduction number, the realised mean over observed periods and the process mean. Counting error falls inversely with the expected count for every target, whereas the process mean also retains between-period transmission variance. Dividing by the previous observed count instead of a known scale raises the counting error by about half under illustrative parameters. When the reporting probability fluctuates between periods, single-period count means and variances cannot separate reporting volatility from transmission variation. We then audit cross-sectional dispersion diagnostics against known truth on the design of 1,748 state-weeks of US influenza hospital admissions (2022--2026). The intercept of an affine fit of relative squared deviation on inverse count is a count-independent relative dispersion: transmission variation under fixed reporting, and its sum with reporting shocks otherwise. In the admissions data it was 0.0150 (95\% confidence interval 0.0101--0.0206). The pooled counting slope was 1.26 (0.81--1.73) times the independent-Poisson reference, and its ratio to a branching-Poisson reference depends on the unknown reporting probability. The dispersion-to-counting crossover count was 37 admissions per week (18--72); tests of equal season intercepts did not reject ($p=0.14$ and $0.21$). Counting slopes and thinning responses do not distinguish overdispersion, under-reporting and heterogeneous transmission; that requires temporal structure or external reporting data.

\medskip
\noindent\textbf{Keywords:} reproduction number; inferential target; surveillance; overdispersion; reporting variation; empirical Bayes
\end{abstract}
